\begin{abstract}
En este proyecto se diseñó un sistema fotovoltaico conectado a la red para reducir la factura eléctrica del edificio de Servicios Complementarios de CoopeAgri R.L., en Pérez Zeledón, que consume \SI{223}{MWh} al año con la tarifa T-CO. A partir de las facturas, los datos solares de \textit{NASA POWER} y \textit{PVGIS} y las fichas técnicas de los equipos, la potencia se dimensionó para cubrir el consumo anual con medición neta sencilla. El sistema propuesto tiene \SI{168.3}{kWp}, con 374 módulos inclinados \SI{10}{\degree} hacia el sur y cinco inversores trifásicos de \SI{25}{kW}. El diseño incluye además la estructura de montaje, el cableado, las protecciones y el diagrama unifilar. Se estima que generaría \SI{219.5}{MWh} al año, el 98.4\,\% del consumo, y que reduciría la factura en ¢14.53 millones anuales (66.4\,\%). Con un precio estimado de ¢119.0 millones, pagando de contado y sin aumento de tarifas, la inversión se recuperaría en 9.2 años, con una TIR de 9.7\,\%.
\end{abstract}

\begin{IEEEkeywords}
sistema fotovoltaico, conexión a red, dimensionamiento, inversor, evaluación económica
\end{IEEEkeywords}
